% Trabajo Práctico Grupal · 2C 2026
% Laboratorio de Métodos Cuantitativos Aplicados a la Gestión · FCE UBA
% Compilar con pdflatex (dos pasadas). Las fechas y el link del formulario
% están en las macros de abajo: se cambian ahí y en ningún otro lado.

\documentclass[11pt,a4paper]{article}

% --------------------------------------------------------------------
% Datos que cambian cada cuatrimestre
% --------------------------------------------------------------------
\newcommand{\Cuatrimestre}{2do Cuatrimestre 2026}
\newcommand{\CuatriCorto}{2C 2026}
\newcommand{\FechaEntrega}{viernes 6 de noviembre de 2026, antes del inicio de la clase}
\newcommand{\FechaDefensa}{viernes 6 de noviembre}
\newcommand{\FechaFormulario}{domingo 18 de octubre de 2026}
\newcommand{\LinkFormulario}{https://docs.google.com/forms/d/e/1FAIpQLScICVFXNGVzvuIHZRwseWbs7Jhze_a4pDw_mYV0V805u7bXLQ/viewform?usp=dialog}

% --------------------------------------------------------------------
\usepackage[utf8]{inputenc}
\usepackage[T1]{fontenc}
\usepackage[spanish,es-noshorthands]{babel}
\usepackage{lmodern}
\usepackage[margin=2.5cm]{geometry}
\usepackage{xcolor}
\usepackage{enumitem}
\usepackage{array}
\usepackage{booktabs}
\usepackage{tabularx}
\usepackage{multicol}
\usepackage{fancyhdr}
\usepackage{titlesec}
\usepackage[most]{tcolorbox}
\usepackage[hidelinks]{hyperref}

\definecolor{azul}{HTML}{243B5E}
\definecolor{enlace}{HTML}{1F4E9A}
\hypersetup{colorlinks=true, urlcolor=enlace, linkcolor=enlace}

\setlength{\parindent}{0pt}
\setlength{\parskip}{0.55em}

\pagestyle{fancy}
\fancyhf{}
\fancyhead[L]{\small\color{gray}Laboratorio de Métodos Cuantitativos \textperiodcentered{} FCE UBA}
\fancyhead[R]{\small\color{gray}Trabajo Práctico Grupal \textperiodcentered{} \CuatriCorto}
\fancyfoot[C]{\small\thepage}
\renewcommand{\headrulewidth}{0.3pt}
\fancypagestyle{plain}{\fancyhf{}\fancyfoot[C]{\small\thepage}\renewcommand{\headrulewidth}{0pt}}

\titleformat{\section}{\large\bfseries}{\thesection.}{0.5em}{}[\vspace{2pt}\titlerule]
\titleformat{\subsection}{\normalsize\bfseries}{\thesubsection.}{0.5em}{}
\titlespacing*{\section}{0pt}{1.6em}{1em}

\newtcolorbox{caja}{breakable, enhanced, colback=gray!4, colframe=gray!60,
  boxrule=0.5pt, arc=2pt, left=6pt, right=6pt, top=4pt, bottom=4pt}
\newtcolorbox{nota}{enhanced, colback=white, colframe=white, borderline west={2pt}{0pt}{azul},
  left=10pt, right=4pt, top=2pt, bottom=2pt}
\newtcolorbox{prompt}[1]{breakable, enhanced, colback=gray!4, colframe=gray!60,
  boxrule=0.5pt, arc=2pt, left=6pt, right=6pt, top=4pt, bottom=4pt,
  title={#1}, coltitle=black, colbacktitle=gray!4, fonttitle=\normalsize,
  attach title to upper, after title={\par\smallskip}}

\begin{document}
\thispagestyle{plain}

% ====================================================================
% Portada
% ====================================================================
\begin{center}
  {\scshape Universidad de Buenos Aires \textperiodcentered{} Facultad de Ciencias Económicas}\\[2pt]
  {\scshape Tecnicatura Universitaria en Gestión y Análisis de Datos en Organizaciones}

  \vspace{1.2em}\rule{\linewidth}{0.8pt}\vspace{0.6em}

  {\LARGE\bfseries Trabajo Práctico Grupal}\\[4pt]
  {\itshape Exploración, análisis y presentación de datos reales asistida por IA}

  \vspace{0.6em}\rule{\linewidth}{0.8pt}
\end{center}

\vspace{0.6em}
\begin{tabularx}{\linewidth}{@{}>{\bfseries}l X@{}}
Modalidad:    & grupos de hasta 5 integrantes \\
Código:       & Google Colab (exportado a \texttt{.ipynb} en GitHub) \\
Informe:      & PDF académico (\LaTeX/Overleaf o Word), subido al repositorio \\
Presentación: & herramienta de IA a elección (Canva, PowerPoint Copilot, Gamma, sitio web, etc.) \\
Repositorio:  & GitHub grupal con commits de todos los integrantes \\
Entrega:      & \FechaEntrega \\
\end{tabularx}

\begin{center}\small \Cuatrimestre\end{center}

% ====================================================================
\section*{Descripción general}
% ====================================================================
Este trabajo integra las herramientas vistas en la materia con una base de datos que cada grupo
elige libremente. El análisis se realiza en Python (Google Colab), se documenta en un informe
con estilo académico (Overleaf/\LaTeX) y se comunica mediante una presentación producida con
herramientas de inteligencia artificial. El uso de IA como asistente está permitido y alentado en
todas las etapas; la interpretación de los resultados y el criterio analítico son responsabilidad
exclusiva del grupo.

Se valorará especialmente la \textbf{originalidad} en la elección del tema, la creatividad en el análisis y
la solidez de las conclusiones.

\begin{nota}
\textbf{Condición de aprobación.} Las tres partes son obligatorias e independientes. La ausencia
de cualquiera de ellas, o la entrega de un notebook sin ejecutar, implica descuento automático
con independencia de la calidad del resto del trabajo.
\end{nota}

% ====================================================================
\section{Inscripción del grupo y elección del tema}
% ====================================================================
La formación del grupo y la elección del dataset se registran mediante el siguiente formulario.
\textbf{Solo un integrante por grupo debe completarlo. Fecha límite para completar el formulario:
\FechaFormulario.}

\begin{center}
\href{\LinkFormulario}{\bfseries $\rightarrow$ Formulario de inscripción grupal (completar antes de comenzar)}
\end{center}

El formulario solicita: nombre y número de registro de cada integrante, nombre de usuario de
GitHub de cada uno, link al repositorio del grupo, dataset elegido, fuente y pregunta principal
que el grupo pretende responder.

\begin{nota}
\textbf{Repositorio de GitHub del grupo.} Cada grupo debe crear un repositorio con todos
los integrantes como colaboradores. El repositorio debe contener: el notebook exportado en
formato \texttt{.ipynb}, el PDF del informe académico, el dataset (o instrucciones para obtenerlo)
y cualquier otro archivo relevante. Se revisará el historial de commits de todos los
integrantes para verificar participación activa y evitar la figura del \emph{freerider}: un integrante
sin commits significativos no acredita el trabajo grupal.
\end{nota}

% ====================================================================
\section{Elección del dataset}
% ====================================================================
El grupo elige libremente el dataset. No existe restricción temática: se aceptan datos económicos,
deportivos, culturales, científicos o de cualquier área de interés. Se destaca y valora la
originalidad en la elección del tema.

Algunos ejemplos posibles (ilustrativos, no limitativos):
\begin{itemize}[itemsep=1pt]
  \item Estadísticas de fútbol: goles, pases, transferencias, rendimiento por liga o club.
  \item Música: reproducciones por género, características de audio, evolución de artistas por período.
  \item Pokémon: estadísticas de combate, tipos, generaciones, evoluciones.
  \item Videojuegos: precios, ventas y valoraciones por plataforma y región.
  \item Datos climáticos, de calidad del aire, de movilidad urbana o de salud pública.
  \item Indicadores macroeconómicos, financieros, de comercio exterior o de mercado laboral.
  \item Series de tiempo de criptomonedas, commodities o mercados financieros emergentes.
  \item Datos de transporte público, Airbnb, delitos urbanos o resultados electorales.
\end{itemize}

\textbf{Fuentes sugeridas:}
\begin{multicols}{2}\small
\begin{itemize}[itemsep=1pt,leftmargin=1.2em]
  \item \href{https://www.indec.gob.ar}{INDEC} --- IPC, EMAE, empleo, comercio exterior
  \item \href{https://datos.gob.ar}{datos.gob.ar} --- datos abiertos del Estado Nacional
  \item \href{https://www.bcra.gob.ar}{BCRA} --- series monetarias y tipo de cambio
  \item \href{https://data.worldbank.org}{Banco Mundial} --- indicadores macro por país
  \item \href{https://ourworldindata.org}{Our World in Data} --- series largas de desarrollo
  \item \href{https://estadisticas.cepal.org}{CEPAL} --- estadísticas de América Latina
  \item \href{https://www.argentina.gob.ar/economia}{Ministerio de Economía} --- presupuesto y deuda
  \item \href{https://finance.yahoo.com}{Yahoo Finance} --- activos vía \texttt{yfinance}
  \item \href{https://www.kaggle.com/datasets}{Kaggle Datasets} --- amplia variedad temática
  \item \href{https://fbref.com}{FBref} / \href{https://sofascore.com}{SofaScore} --- fútbol
  \item \href{https://developer.spotify.com/documentation/web-api}{Spotify API} --- música
  \item \href{https://pokeapi.co}{PokéAPI} --- Pokémon
  \item \href{https://www.pro-football-reference.com}{Pro Football Reference} --- NFL
  \item \href{https://www.basketball-reference.com}{Basketball Reference} --- NBA
  \item \href{https://openweathermap.org/api}{OpenWeatherMap API} --- clima
  \item \href{https://climatologia.meteored.com.ar}{Meteorología Argentina} --- clima local
  \item \href{https://www.imdb.com/interfaces/}{IMDb Datasets} --- cine y series
  \item \href{https://data.un.org}{UN Data} --- estadísticas de Naciones Unidas
  \item \href{https://data.oecd.org}{OECD Data} --- países desarrollados
  \item \href{https://opendata.socrata.com}{Socrata Open Data} --- datos urbanos y gubernamentales
  \item \href{https://datos.jus.gob.ar}{datos.jus.gob.ar} --- justicia y seguridad
  \item \href{https://github.com/datasets}{GitHub Datasets} --- datasets curados en repositorios públicos
\end{itemize}
\end{multicols}

% ====================================================================
\section{Parte 1 --- Análisis en Python (Google Colab)}
% ====================================================================
El notebook debe estar \textbf{completamente ejecutado} al momento de la entrega y exportado en
formato \texttt{.ipynb} al repositorio de GitHub. Cada bloque de código debe ir precedido por una celda
Markdown que explique qué se hace y por qué.

\emph{Sobre el uso de IA en el análisis.} Se recomienda construir un \textbf{agente con contexto enriquecido}:
embeber al modelo (Claude, ChatGPT u otro) con el enunciado de este trabajo, los notebooks y
PDFs del repositorio del curso que consideren relevantes, y cualquier material propio del grupo
(diccionario de datos, paper de referencia, etc.). Cuanto más contexto recibe el agente, mejores y
más pertinentes serán sus sugerencias.

Para optimizar el trabajo en equipo y con la IA se recomienda usar \textbf{proyectos} (Claude Projects,
ChatGPT Projects) donde todos los integrantes comparten el mismo contexto, y explorar
repositorios de prompts especializados como
\href{https://github.com/f/awesome-chatgpt-prompts}{\texttt{awesome-chatgpt-prompts}} o
\href{https://github.com/dair-ai/Prompt-Engineering-Guide}{\texttt{Prompt Engineering Guide}}
para encontrar instrucciones ya probadas. A continuación, dos ejemplos de prompts de esos
repositorios adaptados a este trabajo:

\begin{prompt}{\textbf{Ejemplo 1 --- Rol de analista de datos} \emph{(adaptado de awesome-chatgpt-prompts)}}
\ttfamily «Actuá como un analista de datos senior. Te voy a compartir un dataset de [tema
elegido]. Tu tarea es: (1) identificar las variables más relevantes para responder
[pregunta de investigación], (2) sugerir transformaciones útiles, (3) proponer
qué tipo de gráfico comunica mejor cada hallazgo, y (4) señalar limitaciones del
dataset que puedan afectar las conclusiones. Respondé en español y justificá cada
decisión.»
\end{prompt}

\begin{prompt}{\textbf{Ejemplo 2 --- Chain-of-thought para interpretación económica} \emph{(adaptado de Prompt Engineering Guide)}}
\ttfamily «Pensá paso a paso. Dado el siguiente resultado de mi análisis: [pegar output
del notebook], explicá qué significa económicamente, si es consistente con la
teoría vista en clase, qué factores podrían explicar un resultado inesperado y
qué conclusión debería incluir en el informe académico.»
\end{prompt}

En el notebook debe indicarse qué celdas contaron con asistencia de IA y qué instrucción (\emph{prompt})
se utilizó.

\subsection{Carga y exploración inicial}
\begin{caja}
\begin{enumerate}[label=\alph*), leftmargin=1.6em, itemsep=2pt]
  \item Cargar el dataset en un \texttt{DataFrame} de pandas y mostrar las primeras y últimas filas
        con \texttt{.head()} y \texttt{.tail()}.
  \item Describir la estructura: dimensiones, columnas, tipos de dato (\texttt{.dtypes}) y cobertura
        temporal o geográfica.
  \item Analizar la calidad: valores nulos (\texttt{.isnull().sum()}), duplicados y estadísticas descriptivas
        (\texttt{.describe()}). Comentar los hallazgos antes de continuar.
\end{enumerate}
\end{caja}

\subsection{Pregunta de investigación}
\begin{caja}
Una vez realizada la exploración inicial del dataset, el grupo debe definir con precisión la
pregunta central que guiará el resto del análisis. En el notebook, incluir una celda Markdown
que responda explícitamente:
\begin{enumerate}[label=\alph*), leftmargin=1.6em, itemsep=2pt]
  \item \textbf{¿Cuál es la pregunta?} Formulada de manera clara y acotada; debe ser respondible
        con los datos disponibles.
  \item \textbf{¿Por qué vale la pena responderla?} Motivación del análisis: qué decisión, fenómeno
        o problema ilumina esta pregunta.
  \item \textbf{¿Qué respuesta encontraron?} Al finalizar el análisis, el grupo vuelve a esta celda y
        completa la respuesta con los hallazgos concretos obtenidos.
\end{enumerate}
Esta pregunta es el hilo conductor del trabajo: el informe, los gráficos y la presentación deben
orbitar alrededor de ella y su respuesta.
\end{caja}

\subsection{Transformaciones y resumen estadístico}
\begin{caja}
\begin{enumerate}[label=\alph*), leftmargin=1.6em, itemsep=2pt]
  \item Realizar al menos \textbf{dos transformaciones justificadas} (filtrar filas, convertir tipos,
        crear columnas derivadas, manejar nulos) y explicar la decisión detrás de cada una.
  \item Calcular al menos \textbf{dos estadísticos de resumen agrupados} con \texttt{.groupby()} e
        interpretar los resultados en relación con la pregunta planteada.
\end{enumerate}
\end{caja}

\subsection{Análisis gráfico}
\begin{caja}
Se deben producir \textbf{mínimo tres gráficos de tipos distintos}:
\begin{enumerate}[label=\alph*), leftmargin=1.6em, itemsep=2pt]
  \item \textbf{Evolución temporal o distribución:} serie principal, histograma o boxplot según
        corresponda.
  \item \textbf{Comparación entre categorías:} barras simples o apiladas con ejes etiquetados y
        unidades.
  \item \textbf{Relación entre dos variables:} diagrama de dispersión o líneas superpuestas; señalar
        patrones de interés.
\end{enumerate}
Cada gráfico requiere título, etiquetas en ambos ejes y celda Markdown con interpretación.
No es suficiente describir lo que se ve; se debe explicar qué significa en el contexto del dataset.
Los gráficos se exportan como imágenes (\texttt{plt.savefig()}) para incluirlos en el informe.
\end{caja}

% ====================================================================
\section{Parte 2 --- Informe académico (a elección: \LaTeX{} o Word)}
% ====================================================================
El grupo debe producir un \textbf{informe académico en PDF} con el estilo de un paper o informe de
investigación aplicada. Se puede elegir \textbf{una de las dos opciones} siguientes para su elaboración.

\begin{caja}
\textbf{Opción A --- \LaTeX{} en Overleaf (recomendada).} Se redacta el informe directamente en
Overleaf, o se solicita a un agente de IA (Claude, ChatGPT u otro) que genere el contenido
en un archivo \texttt{.txt} con formato \LaTeX, que luego se compila en Overleaf. Esta opción permite
mayor control tipográfico, numeración automática de figuras y referencias cruzadas. El link al
proyecto Overleaf debe incluirse en el repositorio de GitHub.
\end{caja}

\begin{caja}
\textbf{Opción B --- Word o Google Docs, exportado a PDF.} Se redacta el informe en Microsoft
Word o Google Docs con asistencia de IA, respetando la misma estructura académica. El
archivo final debe exportarse a PDF y subirse al repositorio de GitHub. Se recomienda usar
estilos de título, numeración de figuras y un índice automático para mantener el aspecto
formal.
\end{caja}

En ambos casos, el documento debe ser autónomo: un lector externo debe poder comprender qué
se hizo, qué se encontró y qué se concluye sin abrir el notebook.

\textbf{Proceso de revisión en dos etapas.} Dado que el informe se redacta con asistencia de IA, es
indispensable un proceso de auditoría humana:
\begin{enumerate}[itemsep=2pt]
  \item \textbf{Revisión intragrupal.} Cada integrante lee el informe de forma independiente y contrasta
        las afirmaciones con el notebook y los datos originales. El grupo consolida los comentarios y
        corrige antes de cerrar el documento. Esta etapa es obligatoria.
  \item \textbf{Validación externa} (recomendada para reducir sesgo). El grupo intercambia el informe
        con otro grupo. Un lector externo detecta errores que la familiaridad con el propio trabajo
        vuelve invisibles. Los comentarios recibidos deben incorporarse o responderse en la reflexión
        metodológica.
\end{enumerate}

\textbf{Estructura mínima esperada:}
\begin{itemize}[itemsep=1pt]
  \item \textbf{§1. Resumen} (\emph{abstract}, 5--8 líneas): síntesis del trabajo, metodología y principales hallazgos.
  \item \textbf{§2. Introducción} (1/2 pág.): motivación, dataset, fuente y pregunta principal.
  \item \textbf{§3. Datos y metodología} (1/2 pág.): variables, estructura, transformaciones y herramientas.
  \item \textbf{§4. Resultados} (1--2 pág.): gráficos exportados desde Colab, tablas e interpretación.
  \item \textbf{§5. Conclusiones} (1/2 pág.): hallazgos principales, limitaciones y posibles extensiones.
  \item \textbf{§6. Reflexión metodológica} (1/4 pág.): uso de IA, revisiones realizadas y aprendizajes.
\end{itemize}

\begin{nota}
\textbf{Extensión.} Entre 3 y 5 páginas, sin contar portada ni referencias. Redacción en tercera
persona o impersonal. El PDF final debe subirse al repositorio de GitHub del grupo.
\end{nota}

% ====================================================================
\section{Parte 3 --- Defensa y presentación con herramientas de IA}
% ====================================================================
Todos los grupos defienden su trabajo en vivo el \textbf{\FechaDefensa}, en el horario
habitual de clase (virtual). Cada grupo dispone de un \textbf{máximo de 10 minutos} (incluyendo preguntas),
por lo que la exposición en sí no debe extenderse más de 5 a 7 minutos. Deben estar todos los
integrantes conectados. Se evalúa la claridad comunicacional y la coherencia con el análisis; no la
complejidad técnica de la producción.

La defensa se apoya en una \textbf{presentación de diapositivas obligatoria} producida con IA, a la
que se puede sumar de forma opcional un video o un sitio web como complemento. La estructura
es la siguiente:

\begin{caja}
\textbf{Presentación de diapositivas --- obligatoria para todos los grupos}

Producida con Canva (Magic Design / Canva IA) o PowerPoint con Copilot. El contenido,
guion y narrativa son responsabilidad del grupo; la IA asiste en diseño, estructura y redacción.
Entre 6 y 10 diapositivas calibradas para una exposición de 5--7 minutos. El PDF o link de la
presentación debe subirse al repositorio de GitHub.
\end{caja}

\textbf{Complemento opcional --- elegir 1 de 2 (o ninguno)}

\begin{caja}
\textbf{Opción 1 --- Video explicativo con IA.} Video de 2--4 minutos que recorre los hallazgos
del análisis. Guion redactado por el grupo y revisado con IA. Narración propia o generada
con ElevenLabs. Edición con CapCut, Canva Video o Descript. El link debe subirse al
repositorio. Durante la defensa, el grupo puede reproducir fragmentos del video como parte
de su exposición.

\smallskip
\emph{Material de referencia:} tutoriales y ejemplos de videos en
\href{https://drive.google.com/drive/folders/1x0gSDeycByC-HAt3AUxWFzJ3Co_mkMpl?usp=drive_link}{drive.google.com $\rightarrow$ Carpeta de recursos: videos y páginas web}
\end{caja}

\begin{caja}
\textbf{Opción 2 --- Sitio web desplegado con IA.} Generado y desplegado con asistencia de
IA: Claude o ChatGPT para el código HTML/CSS/JS; Streamlit con deploy en Streamlit
Community Cloud; o GitHub Pages. El sitio debe incluir los tres gráficos, los resultados y
las conclusiones. El link público debe estar en el repositorio. Durante la defensa, el grupo
puede navegar el sitio en vivo como complemento a las diapositivas.

\smallskip
\emph{Material de referencia:}
\begin{itemize}[itemsep=1pt]
  \item Guía de introducción a HTML, CSS y JS:
        \href{https://github.com/Datso653/Laboratorio-de-metodos-Cuantitativos-Aplicados-a-la-gestion/blob/main/PPT's/Complemento\%20TP/HTML\%2CCSS\%20\%26\%20JS\%20-\%20Desarrollo\%20pag\%20web.pdf}{GitHub del curso $\rightarrow$ HTML, CSS \& JS -- Desarrollo pág.\ web}
  \item Tutoriales y ejemplos adicionales:
        \href{https://drive.google.com/drive/folders/1x0gSDeycByC-HAt3AUxWFzJ3Co_mkMpl?usp=drive_link}{drive.google.com $\rightarrow$ Carpeta de recursos: videos y páginas web}
\end{itemize}
\end{caja}

\begin{nota}
\textbf{Nota.} Si el grupo encuentra otra forma creativa de articular IA con una presentación efectiva
(por ejemplo, Gamma.app o Napkin.ai), puede proponerla al docente para su aprobación.
El tiempo de cada grupo está sujeto a la cantidad total de grupos inscriptos; organicen el
contenido para no superar los 5--7 minutos de exposición.
\end{nota}

% ====================================================================
\section*{Entrega --- \FechaEntrega}
% ====================================================================
El entregable es un \textbf{único documento PDF} enviado a través del campus virtual, que debe
contener:
\begin{enumerate}[itemsep=2pt]
  \item \textbf{Portada del grupo:} nombre completo, número de registro y nombre de usuario de
        GitHub de cada integrante, y link al repositorio de GitHub del grupo.
  \item \textbf{Entregable de la Parte 3:} link a la presentación, al sitio público y/o al video según la
        opción elegida.
\end{enumerate}

El repositorio de GitHub debe contener al momento de la entrega, como mínimo: el notebook
exportado en formato \texttt{.ipynb}, el PDF del informe académico, el dataset o instrucciones para
descargarlo, y el archivo o link de la presentación. Se revisará el historial de commits de
cada integrante.

\begin{nota}
\textbf{Participación individual.} Un integrante sin commits significativos en el repositorio no
acredita el trabajo grupal, independientemente de lo declarado en el PDF de entrega.
\end{nota}

\vfill
\begin{center}\small\color{gray}
Consultas por correo electrónico o durante la clase práctica.\\
Laboratorio de Métodos Cuantitativos Aplicados a la Gestión \textperiodcentered{} FCE UBA \textperiodcentered{} \Cuatrimestre
\end{center}

\end{document}
